\section{Related Work}\label{sec:related}

\subsection{Binary-tree adversaries}

Atallah and Kosaraju~\cite{atallah} gave an efficient sorting adversary with
a lower bound of $\tfrac12 n\log n$ comparisons. Related adversaries occur in the work of
Borodin, Guibas, Lynch, and Yao~\cite{borodin} on searching after partial
sorting, Richards and Vaidya~\cite{richards} on the distribution of
comparisons among keys, and Brodal, Chaudhuri, and
Radhakrishnan~\cite{brodal} on maintaining the minimum.
The binary-tree rule of Borodin et al. and Brodal et al. provides the
starting point for our construction.

Place every item at the root of an ordered binary tree. Items in disjoint
subtrees have the left-to-right order of those subtrees. A comparison changes
the placement as follows.
\begin{itemize}
\item If both items occupy the same node, move one to its left child and
the other to its right child, and answer accordingly.
\item If one item is at an ancestor of the other, move the higher item to
the child whose subtree does not contain the lower item.
\item If their nodes are unrelated by ancestry, answer from their
left-to-right order, without moving either item.
\end{itemize}
For a concrete example, start with $A,B,C,D$ at the root.
Answering $A<B$ moves $A$ to the left child and $B$ to the right, while
$C,D$ remain at the root. If the sorter next compares $A$ and $C$, the
ancestor rule moves $C$ to the right child and gives $A<C$.
Comparing $D$ and $B$ then moves $D$ to the left child and gives $D<B$.
The two children now contain $\{A,D\}$ and $\{B,C\}$.
A subsequent comparison of $D$ with $C$ is answered $D<C$ without moving
either item. The orders within $\{A,D\}$ and within $\{B,C\}$ remain
unresolved. Figure~\ref{fig:binary-example} shows the placements.

\begin{figure}[htbp]
\centering
\begin{tikzpicture}[>=Latex,font=\small,
  group/.style={draw=ink!65,rounded corners=2pt,fill=pale,
                minimum height=7mm,minimum width=17mm}]
\begin{scope}[shift={(-3.4,0)}]
\node at (0,1.25) {After answering $A<B$};
\node[group] (root1) at (0,.5) {$C,\ D$};
\node[group] (left1) at (-1.3,-.75) {$A$};
\node[group] (right1) at (1.3,-.75) {$B$};
\draw[->] (root1)--(left1);
\draw[->] (root1)--(right1);
\end{scope}
\draw[->,thick,accent] (-.9,-.1)--(.9,-.1);
\begin{scope}[shift={(3.4,0)}]
\node at (0,1.25) {After also answering $A<C$ and $D<B$};
\node[circle,draw=ink,fill=white,minimum size=4mm] (root2) at (0,.5) {};
\node[group] (left2) at (-1.3,-.75) {$A,\ D$};
\node[group] (right2) at (1.3,-.75) {$B,\ C$};
\draw[->] (root2)--(left2);
\draw[->] (root2)--(right2);
\end{scope}
\end{tikzpicture}
\caption{The binary-tree adversary. The first comparison separates two
items at the same node. The next two comparisons each move an item from
an ancestor into the opposite child. In the final state, comparing $D$
with $C$ requires no movement.}
\label{fig:binary-example}
\end{figure}

Each comparison increases the sum of item depths by at most two.
At completion, occupied nodes must be distinct and unrelated by ancestry:
otherwise two represented items can occur in either relative order, leaving
an unresolved comparison. Associate the root with $(0,1)$, and associate
each child with the corresponding half of its parent's interval.
A node of depth $d_i$ therefore represents an interval of length $2^{-d_i}$.
At completion these intervals are disjoint, giving
\[
\sum_i2^{-d_i}\le1,
\qquad
2^{-\frac1n\sum_i d_i}\le\frac1n\sum_i2^{-d_i}\le\frac1n.
\]
Convexity gives the bound on $2^{-\frac1n\sum_i d_i}$. Hence $\sum_i d_i\ge n\log n$, and
at least $\tfrac12 n\log n$ comparisons are
necessary~\cite[Section~4.1]{brodal}.

The tree representation can settle relations that have never been queried.
For example, start again with four items at the root and answer $A<B$ and
$C<D$. Placing $A,C$ in the left child and $B,D$ in the right child also
commits the adversary to $A<D$ and $C<B$. Neither relation follows from
those two answers alone.

\subsection{The adversary of Kaplan, Zamir, and Zwick}

Kaplan, Zamir, and Zwick~\cite{kzz} replace each step of a four-child tree
by a small directed acyclic graph (DAG), allowing an item to move towards a group of consecutive
children before selecting one child. Number the children $0,1,2,3$.
The intermediate node $[i,j]$ represents the possible destinations
$i,i+1,\ldots,j$. The parent itself is $[0,3]$, and $[i,i]$ is child $i$.
For $i<j$, a downward step removes one endpoint of the range, reaching
$[i,j-1]$ or $[i+1,j]$. Figure~\ref{fig:kzz-lattice} shows this interval
lattice. For example, $[0,2]$ and $[1,3]$ both lead to $[1,2]$; the
intermediate nodes have shared descendants. The same construction is
repeated at every node of the four-child tree.

\begin{figure}[htbp]
\centering
\begin{tikzpicture}[>=Latex,font=\small,
  state/.style={draw=ink!65,rounded corners=2pt,fill=pale,
                minimum height=7mm,minimum width=15mm},
  leaf/.style={state,fill=white}]
\node[state] (n03) at (0,0) {$[0,3]$};
\node[state] (n02) at (-1.25,-1) {$[0,2]$};
\node[state] (n13) at (1.25,-1) {$[1,3]$};
\node[state] (n01) at (-2.5,-2) {$[0,1]$};
\node[state] (n12) at (0,-2) {$[1,2]$};
\node[state] (n23) at (2.5,-2) {$[2,3]$};
\node[leaf] (n0) at (-3.75,-3) {$0$};
\node[leaf] (n1) at (-1.25,-3) {$1$};
\node[leaf] (n2) at (1.25,-3) {$2$};
\node[leaf] (n3) at (3.75,-3) {$3$};
\foreach \a/\b in {n03/n02,n03/n13,n02/n01,n02/n12,n13/n12,n13/n23,n01/n0,n01/n1,n12/n1,n12/n2,n23/n2,n23/n3}
  \draw[->] (\a)--(\b);
\end{tikzpicture}
\caption{One step of the four-child tree in the adversary of Kaplan,
Zamir, and Zwick~\cite[Section~5]{kzz}. An intermediate node records a
range of possible children. In particular, an item at $[1,3]$ has not
committed to child $1$, $2$, or $3$.}
\label{fig:kzz-lattice}
\end{figure}

Within a node, the adversary stores unpaired items and disjoint ordered
pairs. There can be many items at a node, but its retained comparison
components have at most two vertices. Tree nodes and intermediate nodes
representing three children may store pairs; intermediate nodes
representing two children store only unpaired items. Comparing two
unpaired items $A,B$ at a tree node can therefore record $A<B$ as a pair
without moving either item.

A comparison between pairs in different nodes illustrates the role of the
lattice. Suppose $A<B$ is retained at $[0,2]$ and $C<D$ at $[1,3]$.
On comparing $A$ and $C$, one permitted response is $A<C$: move $A$ into
child $0$ and $B$ into $[1,2]$, leaving the pair $C<D$ at $[1,3]$.
Figure~\ref{fig:kzz-pairs} shows the change. The old answer $A<B$ is now
represented by the items' positions, while $C<D$ remains a retained
pair~\cite[Table~11]{kzz}.

\begin{center}
\begin{minipage}{\linewidth}
\captionsetup{hypcap=false}
\centering
\begin{tikzpicture}[>=Latex,font=\small,
  state/.style={draw=ink!35,rounded corners=2pt,fill=white,
                minimum height=7mm,minimum width=12mm},
  occupied/.style={state,draw=ink!70,fill=pale},
  item/.style={circle,draw=ink,fill=white,minimum size=5mm,inner sep=1pt},
  queried/.style={item,fill=accent!18,thick}]
\begin{scope}[shift={(-3.6,0)}]
\node at (0,.7) {Before comparing $A$ and $C$};
\node[state] (l03) at (0,0) {$[0,3]$};
\node[occupied,minimum width=23mm,minimum height=11mm] (l02) at (-1.35,-1.25) {};
\node[font=\footnotesize] at (-1.35,-0.94) {$[0,2]$};
\node[occupied,minimum width=23mm,minimum height=11mm] (l13) at (1.35,-1.25) {};
\node[font=\footnotesize] at (1.35,-0.94) {$[1,3]$};
\node[state] (l01) at (-2.7,-2.65) {$[0,1]$};
\node[state,minimum height=11mm] (l12) at (0,-2.65) {};
\node[font=\footnotesize] at (0,-2.34) {$[1,2]$};
\node[state] (l23) at (2.7,-2.65) {$[2,3]$};
\node[state,minimum width=11mm,minimum height=11mm] (l0) at (-2.85,-4) {};
\node[font=\footnotesize] at (-2.85,-3.69) {$0$};
\node[state,minimum width=11mm,minimum height=11mm] (l1) at (-0.95,-4) {};
\node[font=\footnotesize] at (-0.95,-3.69) {$1$};
\node[state,minimum width=11mm,minimum height=11mm] (l2) at (0.95,-4) {};
\node[font=\footnotesize] at (0.95,-3.69) {$2$};
\node[state,minimum width=11mm,minimum height=11mm] (l3) at (2.85,-4) {};
\node[font=\footnotesize] at (2.85,-3.69) {$3$};
\draw[ink!35] (l03)--(l02);
\draw[ink!35] (l03)--(l13);
\draw[ink!35] (l02)--(l01);
\draw[ink!35] (l02)--(l12);
\draw[ink!35] (l13)--(l12);
\draw[ink!35] (l13)--(l23);
\draw[ink!35] (l01)--(l0);
\draw[ink!35] (l01)--(l1);
\draw[ink!35] (l12)--(l1);
\draw[ink!35] (l12)--(l2);
\draw[ink!35] (l23)--(l2);
\draw[ink!35] (l23)--(l3);
\node[queried] (lA) at (-1.8,-1.44) {$A$};
\node[item] (lB) at (-0.9,-1.44) {$B$};
\node[queried] (lC) at (0.9,-1.44) {$C$};
\node[item] (lD) at (1.8,-1.44) {$D$};
\draw[->,thick,accent] (lA)--(lB);
\draw[->,thick,accent] (lC)--(lD);
\end{scope}
\begin{scope}[shift={(3.6,0)}]
\node at (0,.7) {After answering $A<C$};
\node[state] (r03) at (0,0) {$[0,3]$};
\node[state,minimum width=23mm,minimum height=11mm] (r02) at (-1.35,-1.25) {};
\node[font=\footnotesize] at (-1.35,-0.94) {$[0,2]$};
\node[occupied,minimum width=23mm,minimum height=11mm] (r13) at (1.35,-1.25) {};
\node[font=\footnotesize] at (1.35,-0.94) {$[1,3]$};
\node[state] (r01) at (-2.7,-2.65) {$[0,1]$};
\node[occupied,minimum height=11mm] (r12) at (0,-2.65) {};
\node[font=\footnotesize] at (0,-2.34) {$[1,2]$};
\node[state] (r23) at (2.7,-2.65) {$[2,3]$};
\node[occupied,minimum width=11mm,minimum height=11mm] (r0) at (-2.85,-4) {};
\node[font=\footnotesize] at (-2.85,-3.69) {$0$};
\node[state,minimum width=11mm,minimum height=11mm] (r1) at (-0.95,-4) {};
\node[font=\footnotesize] at (-0.95,-3.69) {$1$};
\node[state,minimum width=11mm,minimum height=11mm] (r2) at (0.95,-4) {};
\node[font=\footnotesize] at (0.95,-3.69) {$2$};
\node[state,minimum width=11mm,minimum height=11mm] (r3) at (2.85,-4) {};
\node[font=\footnotesize] at (2.85,-3.69) {$3$};
\draw[ink!35] (r03)--(r02);
\draw[ink!35] (r03)--(r13);
\draw[ink!35] (r02)--(r01);
\draw[ink!35] (r02)--(r12);
\draw[ink!35] (r13)--(r12);
\draw[ink!35] (r13)--(r23);
\draw[ink!35] (r01)--(r0);
\draw[ink!35] (r01)--(r1);
\draw[ink!35] (r12)--(r1);
\draw[ink!35] (r12)--(r2);
\draw[ink!35] (r23)--(r2);
\draw[ink!35] (r23)--(r3);
\node[item] (rA) at (-2.85,-4.17) {$A$};
\node[item] (rB) at (0,-2.82) {$B$};
\node[item] (rC) at (0.9,-1.44) {$C$};
\node[item] (rD) at (1.8,-1.44) {$D$};
\draw[->,thick,accent] (rC)--(rD);
\end{scope}
\end{tikzpicture}
\captionof{figure}{A comparison between pairs at different nodes in the
adversary of Kaplan, Zamir, and Zwick~\cite[Table~11]{kzz}. Shaded items
are queried. Blue arcs are retained comparisons; gray lines are lattice
edges. The new answer $A<C$ is represented by placement in disjoint child
ranges, so it does not become an arc joining the two pairs.}
\label{fig:kzz-pairs}
\end{minipage}
\end{center}

A retained pair also restricts its members' possible destinations: they
cannot enter the same child of the underlying four-child tree. For
example, if $A<B$ is stored at $[0,2]$, then $A$ can enter only child $0$
or $1$, and $B$ can enter only child $1$ or $2$. These restrictions are
used when choosing subsequent comparison answers. They distinguish the
pair state from simply recording one inequality between otherwise
unrestricted items~\cite[Section~5]{kzz}.

Their analysis measures progress using item depths, intermediate
placements, and retained pairs. The comparison rules bound the increase
in this measure and give a lower bound of
$\frac{37}{64}n\log n$ comparisons, improving the earlier
$\frac12 n\log n$ bound. Their concluding discussion proposes allowing
more complicated comparison components than pairs, and asks whether an
efficient deterministic adversary can force
$(1-o(1))n\log n$ comparisons~\cite[Section~6]{kzz}.
